\section{Behavioral-study implementation and evaluation details}
\label{sec:method-contract}
This appendix explains how the behavioral studies construct and use
executable checks (Appendix~\ref{sec:executable-analysis}). Section~\ref{sec:method}
describes the strategy memory used in planning and error recovery.
\subsection{Constructing historical examples}
A language-model writer reads the historical issue, patch, and relevant code
to propose a requirement and runnable example. Execution on the historical
versions and researcher review check that the example exercises the intended
behavior. Accepted examples are stored with their requirement, input roles,
program, and before/after execution results.

\subsection{Retrieval and repair settings}
The agent uses its full current context to select applicable source records
from the experience pool. Search and graph links help locate requirements
and their validated examples; complete records preserve input types and
argument roles for the subsequent checks.

Within each comparison, agents share the target code, backbone, tools,
instructions, and repair budget. The comparison changes the supplied
memory or feedback, as specified in Appendix~\ref{sec:development-details}.

\subsection{Preparing the current executable program}
\label{sec:public-program}
The development range example extracts model and test snippets from the
public issue and adds researcher-written setup and observation code. The
range and timezone adapters are also researcher-written. The case inputs,
expected behaviors, and control conditions are given in
Appendix~\ref{sec:development-details}.

\subsection{Composition interface}
The composer is the procedure that builds a combined check. It takes the
current program $P_q$, memory $m$, input bindings, and available adapters.
A binding identifies which object plays each role, together with its type,
constructor, argument position, and code location. The composer returns the
check, code that inspects its outcome (the observer), the mapping used, and
whether construction succeeded.
The steps are:
\begin{enumerate}
\item Match an adapter to the current operation and its type requirements.
\item Map current objects to historical input roles, preserving both triggers
and the current operation sequence.
\item Construct the program and check its fixture, bindings, and observer.
The expected outcome comes from the requirement, independently of the candidate patch.
\item Execute the check and return a behavioral failure as repair feedback.
Record setup or observation errors separately so they can be distinguished
from failures of the requirement.
\end{enumerate}

The current adapters cover two Django operation families. The range adapter
transforms the public query's syntax tree, changing an equality lookup for
$x$ into a range with both endpoints equal to $x$, using the historical
named-tuple type $R$:
\begin{equation}
 P_q(x)\longmapsto P_q^{\mathrm{range}}(R(x,x)).
 \label{eq:composition}
\end{equation}
When equality and this one-point range select the same normalized, non-null
value, the expected rows stay unchanged. The adapter preserves the lazy User
and nested query; replacing that object with a lazy integer would lose the
interaction. An ordinary-tuple range checks preservation of existing behavior.
The timezone adapter constructs combinations of database, active application,
and destination roles. Extraction code recovers the bindings these adapters
support. Adding an operation family requires another adapter and its
development tests.
